% GENERATED FILE. Edit canonical lessons, then run npm run generate:books.
% canonical-source-hash: fnv1a64:501d33703c5ffa29
% canonical-lessons: AR-C18-al-taqs

\chapter{Weather}
\label{ch:ar-weather}

This chapter is generated from the canonical micro-lessons used by Language Ladder.
Each section stays independently resumable and preserves the authored prerequisite order.

\section[\ar{الطقس}]{\ar{الطقس} — weather and ritual, one root}

\label{lesson:AR-C18-al-taqs}



\begin{quote}
\textbf{Warm-up.} {[}PAUSE 2s{]} Here's an unexpected pairing: Arabic's word for "weather" shares its root with the word for religious "rites."
\end{quote}

\subsection*{You'll want to know — \ar{طقس} — weather, and its plural means "rites"}
\textbf{\ar{الطقس}} (\emph{aṭ-ṭaqs}) = "\textbf{the weather}" — from root \textbf{\ar{ط}-\ar{ق}-\ar{س}} (\emph{ṭ-q-s}). Its plural, \textbf{\ar{طقوس}} (\emph{ṭuqūs}), means "\textbf{rites, religious ceremonies, liturgy}" — the same root covering both "atmospheric conditions" and "prescribed ritual observances." (Casual speech often prefers \textbf{\ar{الجو}} (\emph{al-jaw}, "the atmosphere") instead, keeping \emph{aṭ-ṭaqs} for more formal or written use.)

\subsection*{You'll want to know — \ar{الجو} \ar{حار} / \ar{بارد} — hot and cold, still verbless}
\begin{quote}
\textbf{\ar{الجو} \ar{حار}.} — "It's hot." — literally "\textbf{the atmosphere {[}is{]} hot}."
\end{quote}

(\emph{al-jaw ḥārr.})

\begin{quote}
\textbf{\ar{الجو} \ar{بارد}.} — "It's cold." — literally "\textbf{the atmosphere {[}is{]} cold}."
\end{quote}

(\emph{al-jaw bārid.})

Same pattern as \textbf{\ar{كم} \ar{عمرك؟}} from the AGE lesson: \textbf{no verb "to be"} at all — just two nouns/adjectives sitting side by side.

\subsection*{You'll want to know — \ar{إنها} \ar{تمطر} — but rain gets its OWN personal verb}
\begin{quote}
\textbf{\ar{إنها} \ar{تمطر}.} — "It's raining." — literally "\textbf{she is raining}." (\emph{innahā}
\end{quote}

tumṭir\emph{ — from \textbf{\ar{مطر}}, }maṭar\emph{, "rain," the same root gives \textbf{\ar{تمطر}}, }tumṭir\emph{, "is raining.")}

Unlike \emph{hot} and \emph{cold}, rain isn't a verbless adjective sentence — Arabic uses an actual \textbf{personal verb}, \emph{tumṭir} ("is raining"), with a \textbf{feminine dummy pronoun}, \textbf{\ar{إنها}} (\emph{innahā}, "she/it"), standing in as the subject. So within Arabic's own weather vocabulary, hot/cold stay verbless while rain gets a real, conjugated verb — echoing the same "rain behaves differently from other weather" pattern you saw in Latin's dedicated \emph{pluit}.

\subsection*{Guided Practice}
{[}PAUSE 1s{]}

\begin{itemize}
\raggedright
  \item {[}YOU SAY: "aṭ-ṭaqs" — weather, same root as "ṭuqūs," rites{]}
  \item {[}YOU SAY: "al-jaw ḥārr. al-jaw bārid." — it's hot/cold, no verb{]}
  \item {[}YOU SAY: "innahā tumṭir" — it's raining, with an actual verb{]}
\end{itemize}

\begin{tcolorbox}[breakable,colback=teal!4,colframe=teal!35!black,title={Wrap-up recall}]
{[}PAUSE 3s{]} What does the plural of \textbf{\ar{طقس}} mean, and what's the connection to "weather"? (\textbf{\ar{طقوس}}, "\textbf{rites, religious ceremonies}" — sharing the same root.) Do \textbf{\ar{الجو} \ar{حار}} and \textbf{\ar{الجو} \ar{بارد}} use a verb? (\textbf{No} — verbless, like \emph{kam ʿumruka?}.) Does \textbf{\ar{إنها} \ar{تمطر}} use a verb? (\textbf{Yes} — \emph{tumṭir}, "is raining," with a feminine dummy pronoun.)
\end{tcolorbox}
